\section{A finite test for Chowla's central nonvanishing question}

Let $\chi$ be a primitive nonprincipal Dirichlet character of conductor
$q$, put $b=(1-\chi(-1))/2$, and write
\[
 \Lambda_\chi(s)=(q/\pi)^{(s+b)/2}\Gamma((s+b)/2)L(s,\chi).
\]
Chowla's question in the quadratic case asks whether every such real
character satisfies $L(1/2,\chi)\ne0$. The repository's one-sided
remainder estimate permits central zeros. Its use of logarithmic derivatives
to convert arithmetic information into zero multiplicities suggests keeping
the residues themselves and using a positive kernel concentrated at a
chosen ordinate. The following kernel uses only values in the half-plane
of absolute convergence. Throughout this section we use Dirichlet GRH.

\paragraph{The finite criterion.}
Fix a real $t$, let $m_t=\operatorname{ord}_{1/2+it}L(s,\chi)$, and put
\[
 H_j(t)=\sqrt j\,\Re\frac{\Lambda_\chi'}{\Lambda_\chi}
       (1/2+\sqrt j+it),\qquad
 B_k(t)=\sum_{j=1}^k(-1)^{j-1}\binom{k}{j} H_j(t),\qquad
 K_k(u)=\prod_{j=1}^k\frac{j}{j+u}.
\]
For every $k\ge1$, counting ordinates $\gamma$ with multiplicity, we have
\begin{equation*}
 B_k(t)=\sum_\gamma K_k((\gamma-t)^2)\ \downarrow\ m_t.
\end{equation*}
Consequently $m_t\le r$ if and only if $B_k(t)<r+1$ for some finite $k$,
for each integer $r\ge0$. For a real quadratic character, the functional
equation has sign $+1$, so $m_0$ is even. In that case $B_k(0)<2$ certifies
central nonvanishing. More generally $B_k(0)<2K_k(g^2)$ for $g\ge0$
certifies the absence of all nontrivial zeros with $|\gamma|\le g$.

\emph{Proof.}
The order-one Hadamard product gives the logarithmic derivative as a
constant plus $\sum_\rho(1/(s-\rho)+1/\rho)$. The sum of
$\Re(1/\rho)$ converges absolutely, by the usual zero count. On the
critical line, away from the zeros, the functional equation
$\Lambda_\chi(s)=\epsilon_\chi\overline{\Lambda_\chi(1-\bar s)}$
makes $\Lambda_\chi'/\Lambda_\chi$ purely imaginary. Since every
$\rho=1/2+i\gamma$, this identifies the real constant in the product and gives
\[
 H_j(t)=\sum_\gamma\frac{j}{j+(\gamma-t)^2}.
\]
Partial fractions give
$\sum_{j=1}^k(-1)^{j-1}\binom{k}{j}j/(j+u)=K_k(u)$.
For $u>0$ the product tends to zero by divergence of $\sum 1/j$.
Dominated convergence applies since the sum for $H_1(t)$ is finite.
Nonreal zeros of the real character occur in pairs of opposite ordinates.
A pair with $|\gamma|\le g$, or a central zero of even order, contributes
at least $2K_k(g^2)$. This proves the claims.

\paragraph{What the convergence rate measures.}
Let $h$ be the least positive distance $|\gamma-t|$ and let $d_h$ count
the zeros at that distance. Then
\[
  B_k(t)-m_t\sim d_h\,\Gamma(1+h^2)k^{-h^2}.
\]
Indeed $K_k(u)=\Gamma(k+1)\Gamma(1+u)/\Gamma(k+1+u)$.
For $u\ge h^2$ the quotient $K_k(u)/K_k(h^2)$ is decreasing with $k$,
is bounded by $K_1(u)/K_1(h^2)$, and tends to zero for $u>h^2$.
The bound is summable over the zeros, so division by $K_k(h^2)$ and
dominated convergence give the result. A zero very close to the target
ordinate therefore slows the convergence by a precise exponent.

\paragraph{The conductor term for a fixed number of samples.}
At $t=0$ and fixed $k$,
\[
 B_k(0)=\tfrac12 c_k\log q+O_k(1),\qquad
 c_k=\sum_{j=1}^k(-1)^{j-1}\binom{k}{j}\sqrt j
 =\frac1{2\sqrt\pi}\int_0^\infty(1-e^{-v})^k\frac{dv}{v^{3/2}}>0 .
\]
The gamma terms are bounded at each of the finitely many arguments, and
the logarithmic derivatives of $L$ are bounded uniformly in $q$ by their
absolutely convergent prime-power series. The integral follows by inserting
$\sqrt j=(2\sqrt\pi)^{-1}\int_0^\infty(1-e^{-jv})v^{-3/2}dv$
into the finite sum. Thus every fixed $k$ eventually gives $B_k(0)>2$
as the conductor grows, regardless of the actual order of vanishing.
This obstruction concerns a fixed depth of this particular test.
The substitution $v=y\log k$ in the integral gives $c_k\sim(\pi\log k)^{-1/2}$.
The factor $(1-k^{-y})^k$ tends to $\mathbf1_{y>1}$ almost everywhere.
On $y\ge1/2$ the other factor $y^{-3/2}$ is integrable. For
$0<v<1/2$ bound $(1-e^{-v})^k$ by $v^k$, and for
$1/2\le v\le (\log k)/2$ bound it by $\exp(-\sqrt k)$.
These bounds make the omitted part negligible after multiplication by
$\sqrt{\log k}$. The surviving integral is $\int_1^\infty y^{-3/2}dy=2$.

\paragraph{A certified case outside Platt's all-moduli search range.}
Let $q=1\,000\,004\,329$ and $\chi(n)=(n/q)$. For this prime conductor
the entire strip $0<\Re s<1$, $|\Im s|\le1/5$ is zero-free for
$L(s,\chi)$. The included Python checker gives the outward enclosures
\[
   0.94936 < B_{10}(0)<1.66669,\qquad
   1.78105 < 2K_{10}(1/25) < 1.78106 .
\]
The modulus was selected by a preliminary numerical search. The proof
uses the final enclosure alone. Platt checked central nonvanishing for
all primitive characters through modulus $2\,000\,000$
(\href{https://doi.org/10.1090/mcom/3077}{Math. Comp. 85 (2016)}).
Platt and Trudgian subsequently studied even quadratic conductors up to $10^9$
(\href{https://arxiv.org/abs/2610.29018}{October 2026 preprint}).
The character here is even and lies just above that bound. We make no
priority claim for this individual conductor.

For reproducibility, at $s=1/2+\sqrt j$ the program evaluates
\[
  \frac{\Lambda_\chi'}{\Lambda_\chi}(s)
  =\tfrac12\bigl(\log(q/\pi)+\psi_0((s+b)/2)\bigr)
    -\sum_n\chi(n)\Lambda(n)n^{-s},
\]
where $\psi_0=\Gamma'/\Gamma$. Trial division verifies primality.
Euler's criterion evaluates the character at primes. All prime powers
through $N=65\,536$ are included. With $\Psi(x)=\sum_{n\le x}\Lambda(n)$
we have $\Psi(2m)-\Psi(m)\le\log\binom{2m}{m}\le 2m\log2$.
Summing at dyadic points and rounding $x$ up to the next such point yields
$\Psi(x)\le4x\log2$. Partial summation bounds the absolute tail by
$4(\log 2)sN^{1-s}/(s-1)$. For $\psi_0$ the recurrence shifts its argument
up by 128, then Binet's formula gives
\[
  \log y-\frac{1}{2y}-\frac{1}{12y^2}
     \le\psi_0(y)\le\log y-\frac{1}{2y}\qquad (y>0).
\]
Indeed the remainder is
$-2\int_0^\infty u\,du/((u^2+y^2)(e^{2\pi u}-1))$.
The checker uses $333/106<\pi<355/113$ and the standard library's
correctly rounded decimal arithmetic at 38 digits with outward expansion
after each operation. The finite calculation and the criterion give one
individual case and a reproducible central zero-free interval.
For background on what is known about the family see Soundararajan,
\href{https://doi.org/10.2307/2261390}{Ann. of Math. 152 (2000)},
and the recent work on pair correlations and Chowla's conjecture,
\href{https://doi.org/10.1007/s00208-026-03383-y}{Math. Ann. (2026)}.
